\subsection{Quadratic signs and the precise arithmetic meaning of Frobenius}

\begin{claim}{C.a1}{Proved}{The two signs are already forced by the matrix norm equations}
Let $\mu\in\widehat{\mathbb Z}^{\times}$, let $t,u,c,d\in\widehat{\mathbb Z}$,
and let $\epsilon_2,\epsilon_3\in\{1,-1\}$ be diagonal signs. If
\[
 \mu(t^2+u^2)=\epsilon_2,\qquad
 \mu(c^2+cd+d^2)=\epsilon_3,
 \tag{C.a1.1}
\]
then
\[
 \epsilon_2=\chi_{-4}(\mu),\qquad
 \epsilon_3=\chi_{-3}(\mu),
\]
where $\chi_{-4}(a)$ is $1$ or $-1$ according as $a\equiv1$ or $-1\pmod4$,
and $\chi_{-3}(a)$ is $1$ or $-1$ according as $a\equiv1$ or $-1\pmod3$.
These are the signs of \S44, folios 365 verso--367, scans 152--155, in
(63), (66)--(68). In particular they do not furnish an additional
arithmetic restriction on a solution for which (C.a1.1) has already been proved.
\end{claim}
\begin{proof}
A square modulo $4$ is $0$ or $1$. Consequently a unit represented by
$t^2+u^2$ is $1\pmod4$. Since both sides of its equation are units,
$\mu\equiv\epsilon_2\pmod4$. Modulo $3$ we have
$c^2+cd+d^2=(c-d)^2$. Its nonzero value is therefore $1$, giving
$\mu\equiv\epsilon_3\pmod3$. Each modulus distinguishes the two diagonal
signs. These are exact calculations in finite rings, also exhaustively
checked in \texttt{code/arithmetic\_checks.py}.
\end{proof}

\begin{claim}{C.a2}{Proved}{Class field interpretation, and the sign $\epsilon'$}
For $\gamma\in G_{\mathbb Q}$ put $\mu=\chi_{\rm cyc}(\gamma)$, the
cyclotomic character for the convention
$\gamma(\zeta_m)=\zeta_m^{\mu\bmod m}$. Then
\[
 \epsilon_2(\gamma)=\frac{\gamma(i)}i,\qquad
 \epsilon_3(\gamma)=\frac{\gamma(\sqrt{-3})}{\sqrt{-3}},\qquad
 \epsilon_2(\gamma)\epsilon_3(\gamma)
       =\frac{\gamma(\sqrt3)}{\sqrt3}.
 \tag{C.a2.1}
\]
Here the quotients are ordinary signs, and changing any of the square roots
does not change its character. If the matrix in scan 155 is written in
the final form of scan 158,
\[
 S=\begin{pmatrix}0&1\\-1&0\end{pmatrix},\qquad
 A=\mu\begin{pmatrix}u&\epsilon_2t\\-t&\epsilon_2u\end{pmatrix},
\]
then $ASA^{-1}=\epsilon_2 S$. Thus the $\epsilon'$ of (67) is
$\epsilon_2$. If $g\rho(g)\rho^2(g)=\epsilon_2$ in the chosen
lift and $A\rho(A)^{-1}=\epsilon''g$, then also
$\epsilon''=\epsilon_2$, as in (68 bis). For an odd prime $p$, its value on arithmetic Frobenius in
$\mathbb Q(i)/\mathbb Q$ is $+1$ exactly when $p\equiv1\pmod4$.
\end{claim}
% anchor:C.a2
\manuscriptanchor{\S 44; folio367; scan155; folio368 verso; scan158}

\begin{proof}
The first equality follows from $i=\zeta_4$ and the fact that the units
modulo $4$ are $1,-1$. Write $\sqrt{-3}=2\zeta_3+1$; exchanging
$\zeta_3$ and $\zeta_3^{-1}$ negates this expression. The second equality
follows, and multiplication gives the third because
$i\sqrt{-3}=-\sqrt3$ after a choice of roots. These are the Artin
characters of the indicated quadratic extensions: at an unramified prime
arithmetic Frobenius sends the residue of a root of unity to its $p$th
power. This proves the Frobenius formula without an unspecified appeal to
class field theory. Geometric Frobenius has the same value because these
characters have order two. Direct multiplication gives $AS=\epsilon_2SA$;
$A$ is invertible by (C.a1.1), so the conjugation formula follows. The product of the three $\rho$-translates
of a coboundary $A\rho(A)^{-1}$ is $1$; the same product for
$\epsilon''g$ is $\epsilon''\epsilon_2$, giving the asserted second sign.

Alternatively, the endomorphism $[i]$ of $y^2=x^3-x$ is
$(x,y)\mapsto(-x,iy)$; Galois conjugation takes it to $[i]$ or $[-i]$
according to its action on $i$. This realizes the same sign intrinsically
on the Tate modules, whenever the matrix represents this complex
multiplication symmetry. The analogous order-three symmetry of
$y^2=x^3-1$ is defined over $\mathbb Q(\zeta_3)$.
\end{proof}
The table, verified by the program, is
\[
\begin{array}{c|rrrr}
 \mu\bmod12&1&5&7&11\\\hline
 \epsilon_2&+1&+1&-1&-1\\
 \epsilon_3&+1&-1&+1&-1\\
 \epsilon_2\epsilon_3&+1&-1&-1&+1
\end{array}
\]
The prime $2$ is ramified in $\mathbb Q(i)$, and the prime $3$ in
$\mathbb Q(\sqrt{-3})$; there is no unramified Frobenius assertion there.
Nor is a global Frobenius multiplier literally the integer $p$ in
$\widehat{\mathbb Z}^{\times}$: $p$ is not a unit at $p$. The statement
uses its residue modulo the conductor, away from $p$.

\begin{claim}{C.a3}{Refuted}{The $2$-adic square criterion in scan 152}
The sentence in \S44, folio 365 verso, scan 152 identifying
$\epsilon_p(\mu)=1$ with $\mu_p\in(\mathbb Z_p^{\times})^2$ cannot
apply to the sign $\epsilon_2$ defined there. The strongest valid uniform
replacement is the explicit pair of conductor-$4$ and conductor-$3$
characters in C.a1; for $p=2$ the square test requires modulus $8$.
\end{claim}
\begin{proof}
For example $5\equiv1\pmod4$, so $\chi_{-4}(5)=1$, whereas the square
of every odd integer is $1\pmod8$; therefore $5\notin
(\mathbb Z_2^{\times})^2$. If an actual profinite unit is wanted, choose
its $2$-component equal to $5$ and every odd component equal to $1$.
The finite counterexample is printed by the program. Conversely an odd
$2$-adic unit is a square exactly when it is $1\pmod8$: necessity is
as above; for sufficiency, starting with a square root modulo $8$, replace
an odd root $r$ modulo $2^k$ by $r+2^{k-1}$ to change its square by
$2^kr$ modulo $2^{k+1}$, and choose the sign of the correction to lift
one step. Compactness gives a $2$-adic root. For an odd prime the usual
lifting argument has invertible derivative $2r$ and proves the square
criterion modulo $p$; at $p=3$ it agrees with C.a1.
\end{proof}
No handwriting is silently changed by this correction. The scan-155 image
supports the displayed modulus $4$; the mathematical error is the overly
uniform square interpretation of a different passage, not that modulus.
The later formula $\mu\equiv\epsilon_2\pmod2$ in scan 165 determines
no sign and must be strengthened to modulus $4$ if it is to do the work
claimed for it.

\subsection{Faithfulness on the modular group and its centralizer}

Throughout this subsection $U=\mathbb P^1_{\mathbb Q}\setminus
\{0,1,\infty\}$,
$F=\pi_1(U_{\overline{\mathbb Q}})\simeq\widehat F_2$,
$H=\widehat{\mathrm{SL}_2(\mathbb Z)}$, and
$P=H/\langle -I\rangle\simeq\widehat{\mathrm{PSL}_2(\mathbb Z)}$.
Every automorphism is continuous and every outer automorphism is taken
modulo conjugation by an element of its target. Compatible geometric
base points identify $F$ with $\ker(P\to\mathfrak S_3)$, up to
$P$-conjugacy. We use the arithmetic actions arising from the stacks
$\mathcal M_{1,1}$ and its rigidification by $\{\pm1\}$.

\begin{claim}{C.a4}{Known}{Exactly what faithfulness modulo geometric monodromy means}
Let $\Delta=\pi_1(E_{\overline{\mathbb Q}}\setminus\{0\})$ for an
elliptic curve $E/\mathbb Q$. Let
\[
 \Psi:\pi_1(\mathcal M_{1,1})\longrightarrow\operatorname{Out}(\Delta),
 \quad N_{\rm ar}=\operatorname{im}\Psi,\quad
 \Sigma=\Psi\bigl(\pi_1(\mathcal M_{1,1,\overline{\mathbb Q}})\bigr).
\]
Then the canonical map $G_{\mathbb Q}\to N_{\rm ar}/\Sigma$ is an
isomorphism of profinite groups. This is faithfulness modulo the geometric
orientation-preserving monodromy image, as asked in \S35--36, folios
254--255 verso, scans 502--505. It is not by itself the assertion
$G_{\mathbb Q}\hookrightarrow\operatorname{Out}(H)$.
\end{claim}
\begin{proof}
The chosen elliptic curve gives a section over $G_{\mathbb Q}$ of the
arithmetic stack extension, and hence identifies its action on $\Delta$
with the usual elliptic-curve outer action. Apply
\citet[Corollary 6.4]{HM2011arith} with $k=\mathbb Q$ and $(g,r)=(1,1)$.
Its precise conclusion is
$\rho_{E\setminus0}^{-1}(\Sigma)\subseteq
\ker(G_{\mathbb Q}\to G_{\mathbb Q})=1$.
Surjectivity of $G_{\mathbb Q}\to N_{\rm ar}/\Sigma$ follows from
surjectivity of the stack extension onto $G_{\mathbb Q}$. A continuous
bijection between compact Hausdorff groups is a topological isomorphism.
The theorem concerns $\Sigma$, whether or not it has already been
identified with $H$; \citet[Corollary 6.5]{HM2011arith} explicitly
separates injectivity of total monodromy from that geometric congruence
question. It also applies with any number field in place of $\mathbb Q$.
\end{proof}

\begin{claim}{C.a5}{Proved}{Faithfulness on $H$ and already on $P$}
The natural continuous outer actions
\[
 G_{\mathbb Q}\longrightarrow\operatorname{Out}(P),\qquad
 G_{\mathbb Q}\longrightarrow\operatorname{Out}(H)
\]
are injective. This proves the stronger statement eventually written in
\S36, folio 270 verso, scan 530, as well as answering the earlier doubt.
\end{claim}
\begin{proof}
The ordered level-two covering of the rigidified modular stack is defined
over $\mathbb Q$. It gives the Galois-compatible exact sequence
$1\to F\to P\to\mathfrak S_3\to1$ of scans 521--523. If
$\gamma$ acts on $P$ by conjugation by $p$, its action on $F$, modulo
$F$-conjugacy, equals the anharmonic permutation represented by
$pF\in\mathfrak S_3$. On $F^{\rm ab}\simeq\widehat{\mathbb Z}^2$,
Galois acts by the scalar $\chi_{\rm cyc}(\gamma)$, whereas
$\mathfrak S_3$ acts by its standard action on the three vectors
$e_0,e_1,e_\infty$ with $e_0+e_1+e_\infty=0$. A scalar permuting these
three distinct vectors is the identity: it cannot send $e_0$ to $e_1$
or $-e_0-e_1$, so fixes $e_0$, and then fixes all three. Thus $p\in F$.
Consequently $\gamma$ acts trivially in $\operatorname{Out}(F)$.
Belyi's faithfulness theorem
\citep[Theorem 4 and its Corollary]{Belyi1980arith}, or the directly
applicable \citet[Corollary 6.3(i)]{HM2011arith}, gives $\gamma=1$.
An inner action on $H$ induces an inner action on $P$, so the second
assertion follows from the first. This argument needs no assertion that
$\operatorname{Out}(H)$ is $\widehat{\mathrm{GT}}$.
\end{proof}

\begin{claim}{C.a6}{Known}{The tripod centralizer used below}
For its standard arithmetic action,
\[
 C_{\operatorname{Out}(F)}(G_{\mathbb Q})=\mathfrak S_3.
\]
\end{claim}
% anchor:C.a6
\manuscriptanchor{\S 36; folios266--270 verso; scans521--530 (cote140-2)}

\begin{proof}
We give the reduction to the primary theorem actually opened, rather
than importing the assertion from LM3. In
\citet[Corollary 6.3(iii),(iv)]{HM2011arith}, take $X=U$, $k=k_0=
\mathbb Q$, and $n=1$. The tripod is of quasi-Belyi type, and the
fibration-preserving qualification $\operatorname{Out}^{F}$ is empty
for one point. The theorem says that
$A=\operatorname{Aut}(U_{\overline{\mathbb Q}}/\mathbb Q)
=G_{\mathbb Q}\times\mathfrak S_3$ embeds and is commensurably terminal
in $\operatorname{Out}(F)$: its commensurator equals $A$.
If $a$ centralizes $G_{\mathbb Q}$, the intersection $A\cap aAa^{-1}$
contains $G_{\mathbb Q}$, of index six in each group. Hence $a\in A$.
It remains to see that $Z(G_{\mathbb Q})=1$. Choose the real algebraic
closure $R\subset\overline{\mathbb Q}\subset\mathbb C$ and its
complex conjugation $c$. An element commuting with $c$ preserves $R$.
Every automorphism of $R$ preserves its order because positive elements
are squares, and fixes its dense rational subfield, hence fixes $R$
pointwise. Thus $C_{G_{\mathbb Q}}(c)=\{1,c\}$. The image of $c$ in
the splitting field of $X^3-2$ is a noncentral transposition of
$\mathfrak S_3$, so $c$ is not central in $G_{\mathbb Q}$. Therefore
$Z(G_{\mathbb Q})=1$, and the centralizer inside $A$ is exactly its
second factor, as asserted.
\end{proof}

\begin{claim}{C.a7}{Proved}{The centralizer on the modular group is a central transvection}
With the arithmetic actions of C.a5,
\[
 C_{\operatorname{Out}(P)}(G_{\mathbb Q})=1,\qquad
 C_{\operatorname{Out}(H)}(G_{\mathbb Q})\simeq C_2.
 \tag{C.a7.1}
\]
The second group is generated by the outer class of
\[
 \theta(h)=(-I)^{\eta(h)}h,
 \qquad \eta:H\longrightarrow C_2,
 \tag{C.a7.2}
\]
where $\eta$ is the unique nonzero character, equivalently the sign of
reduction modulo $2$ in $\mathfrak S_3$.
\end{claim}
% anchor:C.a7
\manuscriptanchor{\S 36; folios266--270 verso; scans521--530 (cote140-2)}

\begin{proof}
We spell out the group-theoretic comparison. Since
$P=C_2\widehat{*}C_3$, every surjection $P\to\mathfrak S_3$ sends the
order-two generator to a transposition and the order-three generator to
a three-cycle. The six choices are conjugate, so their common kernel
$F$ is characteristic. The induced action of $P/F$ on $F^{\rm ab}$ is
faithful, and the free profinite group $F$ is centre-free; thus
$C_P(F)=1$. The centre-free property is the Corollary to Lemma 1 of
\citet[\S1, pp.~248--249]{Belyi1980arith}; its proof uses the universal
lifting property of $\widehat F_2$ and noncommuting central extensions
of elementary abelian finite quotients.
An automorphism of $P$ restricting to the identity on $F$ is the identity:
for $p\in P$, its discrepancy $\alpha(p)p^{-1}$ centralizes $F$.
It follows that restriction gives an injection
\[
 \operatorname{Out}(P)\hookrightarrow
 N_{\operatorname{Out}(F)}(\mathfrak S_3)/\mathfrak S_3.
 \tag{C.a7.3}
\]
If $a\in\operatorname{Out}(P)$ centralizes Galois, choose a restriction
$\widetilde a\in N_{\operatorname{Out}(F)}(\mathfrak S_3)$. For every
$\gamma\in G_{\mathbb Q}$ the commutator
$[\widetilde a,\rho_F(\gamma)]$ belongs to $\mathfrak S_3$.
Both $\rho_F(\gamma)$ and its $\widetilde a$-conjugate centralize
$\mathfrak S_3$, since the anharmonic maps are defined over $\mathbb Q$.
This commutator consequently lies in $Z(\mathfrak S_3)=1$.
C.a6 yields $\widetilde a\in\mathfrak S_3$, and (C.a7.3) gives $a=1$.

Now $Z(P)=1$ by $C_P(F)=1$, so $Z(H)=\langle-I\rangle$ is
characteristic. If an automorphism of $H$ is inner after passage to $P$,
remove the inner automorphism by lifting its conjugator to $H$. The
result has the form $h\mapsto(-I)^{b(h)}h$ for a continuous character
$b:H\to C_2$ vanishing on $-I$. Abelianizing the presentation
$S^4=1$, $S^2=(ST)^3$ gives $H^{\rm ab}=\mathbb Z/12$, with $T$ a
generator and $-I$ corresponding to $6$. Therefore exactly two such
characters occur, zero and $\eta$. Their automorphisms are invertible;
$\theta^2=1$ because $\eta(-I)=0$. The nonzero one is not inner,
since it acts on $H^{\rm ab}$ by multiplication by $7$, rather than $1$.
It commutes with every automorphism of $H$, because both $-I$ and
$\eta$ are characteristic. The kernel of
$\operatorname{Out}(H)\to\operatorname{Out}(P)$ is therefore this
central $C_2$. Combining it with the already proved centralizer on $P$
gives (C.a7.1).
\end{proof}
The program checks (C.a7.2) on every element and a generating pair of
$\mathrm{SL}_2(\mathbb Z/n)$ for $n=2,4,6,8,12,24$, and searches every
inner conjugator. It is trivial at $n=2$ and noninner at each of the other
five levels. This finite check supports the explicit formula; the proof
above is the assertion about the full profinite group. The theorem of
\citet[Theorem B, Theorem 4.3 and Corollary C]{MN2022arith} concerns
four or more strands and $\widehat\Gamma_{1,2}$. It must not be
specialized to three strands or used to identify
$\operatorname{Out}(H)$ with $\widehat{\mathrm{GT}}$.

\subsection{The lifts through $\{\pm1\}$ and the withdrawn nonexistence doubt}

\begin{claim}{C.a8}{Proved}{The obstruction and the torsor of lifts}
Let $1\to A\to E\xrightarrow{q}B\to1$ be a central extension of
profinite groups with $A=C_2$, and let $j:D\to B$ be continuous.
The obstruction to a continuous lift $\widetilde j:D\to E$ is the
pullback class $j^*[E]\in H^2_{\rm cont}(D,A)$. When a lift exists,
its set is a torsor under $H^1_{\rm cont}(D,A)=\operatorname{Hom}_{\rm cont}(D,A)$.
For $D=\widehat F_2$ this obstruction vanishes and there are exactly
four lifts of a fixed $j$. For $D=F_2$ the same statements hold in the
discrete category.
\end{claim}
% anchor:C.a8
\manuscriptanchor{\S 36; folios267--268; scans523--525 (cote140-2)}

\begin{proof}
A finite-kernel surjection of profinite groups admits a continuous
set-theoretic section: take an open subgroup meeting $A$ trivially,
use its inverse section over its open image, and choose representatives
on finitely many clopen cosets. Such a section gives a continuous factor
set $c\in Z^2(B,A)$. A lift is a function $b:D\to A$ with
$\delta b=j^*c$ (up to the immaterial inverse convention for $c$).
This proves the obstruction criterion. Given two lifts, their pointwise
ratio is a continuous homomorphism $D\to A$ because $A$ is central;
every such homomorphism changes a lift to another, freely and
transitively. A lift for $\widehat F_2$ is obtained by choosing
arbitrary lifts of its two free generators and using the universal
property. The choices in each fiber are exactly two. The same argument
works for the discrete free group.
\end{proof}
For specified genuine actions of a group $\Gamma$ on this diagram, the
set of equivariant lifts, if nonempty, is the fixed torsor under
$H^1(D,A)^\Gamma$. A fixed point is an extra condition on an affine
$\Gamma$-action, not an automatic consequence of $H^2(D,A)=0$.
For outer arithmetic actions one must use the full arithmetic extensions;
scan 524 makes precisely this correction. The next result solves that
extra condition in the case actually considered there.

\begin{claim}{C.a9}{Refuted}{The family doubted on scan 524 exists over $\mathbb Q$}
The nonexistence suggested in \S36, folio 267 verso, scan 524 is false.
Indeed it is explicitly withdrawn in the checked continuation at scans
527--528, folio 269. Over
$U=\operatorname{Spec}\mathbb Q[\lambda,1/\lambda,1/(\lambda-1)]$,
\[
 E_\lambda:\quad Y^2Z=X(X-Z)(X-\lambda Z)
 \tag{C.a9.1}
\]
with origin $[0:1:0]$ and ordered nonzero two-torsion points
$(0,0),(1,0),(\lambda,0)$ is the required family. It gives a morphism
$U\to\mathcal M_{1,1}$ lifting the tautological map to its
$\{\pm1\}$-rigidification, and therefore a morphism of arithmetic
fundamental-group extensions over $G_{\mathbb Q}$.
\end{claim}
\begin{proof}
The discriminant is $16\lambda^2(\lambda-1)^2$, a unit on $U$;
the projective cubic is smooth, with the stated origin and three distinct
nonzero points of order two. Its quotient by negation is the projective
$X/Z$-line branched at $0,1,\lambda,\infty$, so its ordered level-two
invariant is exactly $\lambda$. Functoriality of \'etale $\pi_1$ gives
the asserted homomorphism, continuous and well-defined up to conjugation
by the geometric target group. This is the homomorphism (40) on scan
528. The polynomial discriminant is also checked symbolically by the
program. The construction works over $\mathbb Z[1/2]$ with
$\lambda$ and $\lambda-1$ inverted, including characteristic $3$;
it makes no assertion about ordinary full level two in characteristic $2$.
\end{proof}

\begin{claim}{C.a10}{Proved}{All four geometric lifts are arithmetic and geometric}
Fix the quotient map $U\to\mathcal M_{1,1}^{\rm rig}$ just used and
fix the Legendre lift as origin. The isomorphism classes of its lifts
through the $C_2$-gerbe $\mathcal M_{1,1}\to\mathcal M_{1,1}^{\rm rig}$
are in bijection with
\[
 H^1_{\rm et}(U,\mu_2)
 =\mathbb Q^{\times}/\mathbb Q^{\times2}\ \oplus\ (\mathbb Z/2)^2.
 \tag{C.a10.1}
\]
They are realized by
\[
 a\lambda^e(\lambda-1)^fY^2Z=X(X-Z)(X-\lambda Z),
 \quad a\in\mathbb Q^{\times}/\mathbb Q^{\times2},\quad e,f\in\{0,1\}.
 \tag{C.a10.2}
\]
The restriction to geometric fundamental groups gives exactly the four
lifts $F\to H$ of $F\hookrightarrow P$ in C.a8. All four are fixed
under the outer $G_{\mathbb Q}$-action, and all arise from (C.a10.2)
already over $\mathbb Q$.
\end{claim}
% anchor:C.a10
\manuscriptanchor{\S 36; folios267--269 verso; scans523--528 (cote140-2)}

\begin{proof}
A choice of object in a gerbe banded by $\mu_2$ identifies its other
objects with $\mu_2$-torsors by their sheaves of isomorphisms. Here the
Legendre family is that object, and the band is $\mu_2$: after rigidifying
by all labels, the only automorphisms are $\pm1$. Let
$R=\mathbb Q[\lambda,\lambda^{-1},(\lambda-1)^{-1}]$. This is a
localization of a principal ideal domain, so $\operatorname{Pic}(R)=0$;
its units are exactly $a\lambda^m(\lambda-1)^n$, with
$a\in\mathbb Q^{\times}$ and $m,n\in\mathbb Z$. The Kummer sequence
therefore gives (C.a10.1), and descent along $w^2=a\lambda^e
(\lambda-1)^f$ gives precisely the quadratic twists (C.a10.2).
Over $\overline{\mathbb Q}$ the constant $a$ becomes a square, whereas
$\lambda$ and $\lambda-1$ have independent odd valuations at $0$ and
$1$; they give the two independent characters of $F$ into $C_2$.
Their four combinations exhaust C.a8. They are defined over $\mathbb Q$,
so all four have arithmetic representatives and are Galois-fixed.
They remain distinct modulo conjugation: a conjugator identifying two
lifts while fixing their quotient inclusion would project to an element
of $C_P(F)=1$, and hence would be central in $H$ and act trivially.
\end{proof}

\begin{claim}{C.a11}{Proved}{The precise anabelian converse here needs no section conjecture}
For the fixed tautological quotient map in scans 523--526, every lift of
the arithmetic fundamental-group homomorphism through $\{\pm1\}$ is
induced by a geometric lift, up to its stated geometric inner ambiguity;
the same is true of every outer-Galois-compatible geometric lift
$F\to H$. No assertion that every section of
$\pi_1(\mathcal M_{1,1})\to G_{\mathbb Q}$ is geometric is used or
implied.
\end{claim}
% anchor:C.a11
\manuscriptanchor{\S 36; folios267--269; scans523--528 (cote140-2)}

\begin{proof}
For arithmetic fundamental groups, fix the lift furnished by C.a9.
The ratio of any other lift is a continuous $C_2$-character of
$\pi_1(U)$, by C.a8. Finite \'etale torsors are classified by the
fundamental group, so
\[
 H^1_{\rm cont}(\pi_1(U),C_2)=H^1_{\rm et}(U,\mu_2).
\]
The corresponding twist in C.a10 induces that character and hence that
lift. For geometric groups, C.a10 realizes all four possible lifts and
therefore all Galois-compatible ones. The weaker datum of a geometric
lift does not determine the arithmetic lift uniquely: the kernel of
restriction in (C.a10.1) is
$\mathbb Q^{\times}/\mathbb Q^{\times2}$. This explains exactly the
ambiguity that the manuscript starts to detect on scan 524. The
$H^2$ comparison suggested on scan 526 is unnecessary after the explicit
neutralization by the Legendre family; if one uses it, both obstruction
classes are already zero.
\end{proof}
Thus the particular request C(h) is settled more strongly than a reduction
to a section conjecture. A classification of arbitrary sections of a
modular curve is a separate assertion with different input data. It would
be a mistake to transfer an unverified statement about such sections to
this finite-band gerbe problem, or conversely to deduce a section theorem
from the elementary lifting result proved here.

\subsection{A characterization of the arithmetic image with specified additional data}

\begin{claim}{C.a12}{Known}{Pop's theorem gives an exact extendability criterion}
Let $X=\mathbb A^2_{\mathbb Q}$ with coordinates $u,v$, and let
$\mathcal V$ be the category with objects $U$ and a basis of sufficiently
small nonempty principal open subsets of $X$ defined over $\mathbb Q$.
Include all inclusions between these opens and the maps to $U$ defined
by the three rational functions $u,v,v-u$, whenever regular with values
in $U$. Put $\pi_{\mathcal V}(Y)=\pi_1(Y_{\overline{\mathbb Q}})$ in
the category of profinite groups and outer continuous homomorphisms.
Then evaluation at $U$ identifies
\[
 \operatorname{Aut}(\pi_{\mathcal V})\xrightarrow{\ \sim\ }
 \rho_U(G_{\mathbb Q})\subset\operatorname{Out}(F).
 \tag{C.a12.1}
\]
In particular an element of any of Grothendieck's candidate groups
acting on $F$ is arithmetic exactly when its outer action extends to
such a compatible family. This answers the characterization question
of \S45, scans 170,178,185, by specifying data beyond the two types.
\end{claim}
% anchor:C.a12
\manuscriptanchor{\S \S 45--46; fols.384,388,391 verso; cote140-3 scans170,178,185}

\begin{proof}
We verify Hypothesis (H) of \citet[Definition/Remark 2.2]{Pop2018arith}.
The function field $\overline{\mathbb Q}(u,v)$ has transcendence degree
$2$ and is generated by the indicated functions. A compatible birational
automorphism has one common automorphism of the tripod on all three
maps. Its restriction to constants is some $\sigma\in G_{\mathbb Q}$;
after removing $\sigma$, the tripod automorphism is one of the six
anharmonic fractional linear functions $r$. Compatibility forces
\[
 r(v)-r(u)=r(v-u).
\]
Only $r(t)=t$ satisfies this rational identity. For $r=1-t$ its discrepancy
is $-1$; for the four remaining nonidentity functions it is a nonzero
rational function, listed in the exact symbolic output of the program.
Thus the birational automorphism fixes $u,v$ and is exactly $\sigma$ on
constants. Conversely every $\sigma$ fixes every rationally defined
open and the three maps. This proves rigidity and hence Hypothesis (H).
Apply \citet[Theorem 2.7(1)]{Pop2018arith}; characteristic zero makes its
valuation-tame fundamental groups equal to the full groups under
consideration. It gives $G_{\mathbb Q}\simeq
\operatorname{Aut}(\pi_{\mathcal V})$. Evaluation is the usual action
on $F$, faithful by C.a5's tripod input, proving (C.a12.1).
\end{proof}
The extra datum is not a finite list of finite quotients: even though
there are only three designated tripod functions, the neighborhood
category includes complements of arbitrarily many divisors in a
surface. The theorem does not say that checking the two original
modular types or any fixed level in the requested computations suffices.

\begin{claim}{C.a13}{Known}{The combinatorial characterization determines a conjugacy class}
Let $\mathcal A$ be the collection of closed subgroups of
$\widehat{\mathrm{GT}}$ satisfying the AA-property of
\citet[Definition 5.12]{HMT2025arith}, including its COF, RGC and
ArBC prerequisites. The maximal elements of
$\mathcal A$ are exactly the
$\widehat{\mathrm{GT}}$-conjugates of $\rho_U(G_{\mathbb Q})$.
This is a characterization of a conjugacy class of subgroups, not a
canonical distinguished embedding or an elementwise equation on
$(\lambda,f)$.
\end{claim}
% anchor:C.a13
\manuscriptanchor{\S \S 45--46; fols.384,388,391 verso; cote140-3 scans170,178,185}

\begin{proof}
The precise theorem is \citet[Theorem 5.17(ii), also Theorem C(ii)]{HMT2025arith}.
We record its reduction to prevent a stronger reading. The arithmetic
subgroup itself satisfies AA (their Remark 5.12.1). Their Theorem
5.15(iv) puts any AA subgroup in a conjugate of the arithmetic subgroup.
Their Lemma 5.16 rules out a proper containment between the relevant
conjugates of the latter. Thus maximal AA subgroups and exactly those
conjugates coincide. Transport the statement along the chosen
identification $F\simeq\widehat\pi_{0,3}$ and, when appropriate, along
the quotient from Grothendieck's extension group to its outer loop
action. Transport changes the embedded subgroup by the declared
conjugacy ambiguity, and changes nothing about maximality or closure.
\end{proof}
The input from which the theorem reconstructs this collection is the
abstract topological group of the second configuration space of the
tripod, $\pi_1(M_{0,5,\overline{\mathbb Q}})$, together with its
intrinsically reconstructed structure. Definition 5.12 quantifies over
all normal open subgroups and cusp sets and requires nondegenerate
Kummer data and compatible fields of functions. These conditions are
substantially more than the order-two and order-three equations on one
copy of $F$. We opened the stated theorem, the reduction just cited,
and Definition 5.12; we do not claim an independent verification of all
125 pages of the argument.

\begin{claim}{C.a14}{Proved}{The arithmetic conclusions have the following graded limits}
On the subgroup $\lambda=1$, the characters $\epsilon_2$ and
$\epsilon_3$ of C.a1--C.a2 are identically $+1$. Their infinitesimal
conditions are therefore vacuous. The order-two transvection in C.a7,
the finite-band lifting choices in C.a8--C.a11, and the inverse-limit
extendability condition of C.a12 are not recovered by merely taking
rational weight-graded Lie algebras of $F$.
\end{claim}
% anchor:C.a14
\manuscriptanchor{\S 44; fols.365 verso--368 verso; cote140-3 scans152--158; \S 36; fols.267--269 verso; cote140-2 scans523--528}

\begin{proof}
The first assertion is evaluation of the two characters at $1$. Any
map from a uniquely divisible unipotent group over a characteristic-zero
field to a group of exponent two is trivial; equivalently the tangent
space of the constant finite group $C_2$ is zero. Consequently neither
four distinct $C_2$-twists nor the nonzero transvection can be encoded by
that tangent space. The criterion C.a12 includes the full geometric
fundamental groups of every object and all comparison maps, whereas
the stated graded construction only retains the lower-central
successive quotients of $F$. No implication from the latter data to the
former is part of Pop's theorem. This last observation identifies the
missing comparison map; it is not an assertion that no other Lie-theoretic
reconstruction could ever exist.
\end{proof}
